%% Full proofs. Verified numerically in the released artifact (theory stages t15, t18, t21f,
%% t32a, t34a, t35a, t36a). Notation matches the body: M = ceil = (|C|+1)/k is the evidence
%% ceiling, every per-flow e-value lies in {0, M}, and a rejection at step t needs evidence
%% >= 1/alpha_t.

\subsection{Group-level marginal e-validity (\cref{lem:groupval})}

\begin{proof}
For a tested group $G_j$, recall $\mathcal{M}_j$ is the $\sigma$-field generated by the pre-committed
metadata that determines \emph{this} hypothesis: the grouping key, the membership indicators
$\mathbf 1\{i\in G_j\}$, the arity $m_j$, the group's position in the pre-committed stream order (and
hence the spending weight $\gamma_j$ it is offered), and the horizon $T$. It carries no information
about any other group's composition or arity, and --- as the scope note below records --- none about
the realised calibration scores. Fix a true-null group $G_j$ (all constituent flows benign) with
$m_j>0$. By linearity of conditional expectation and $\mathcal{M}_j$-measurability of the membership
indicators and of $m_j$,
\begin{align*}
\mathbb{E}[\Ev(G_j)\mid\mathcal{M}_j]
 &=\mathbb{E}\big[\tfrac1{m_j}\textstyle\sum_{i\in G_j} e_i \,\big|\,\mathcal{M}_j\big]\\
 &=\tfrac1{m_j}\textstyle\sum_{i\in G_j}\mathbb{E}[e_i\mid\mathcal{M}_j]\ \le\ 1,
\end{align*}
the inequality by \cref{assump:groupval}. Taking expectations and applying the tower property gives
$\mathbb{E}[\Ev(G_j)]=\mathbb{E}\big[\mathbb{E}[\Ev(G_j)\mid\mathcal{M}_j]\big]\le1$. No independence or
exchangeability \emph{among} the $\{e_i\}$ is used, so within-group dependence is unrestricted; and
because e-LOND and online e-BH control $\FDR$ under arbitrary dependence given marginally valid
per-hypothesis e-values, arbitrary between-group dependence is permitted as well. An empty group
($m_j=0$) tests no hypothesis; assigning it the trivial $\Ev=1$ is harmless. For the two-point
conformal e-value $e_i=\ceil\,\mathbf 1\{K_i\le k\}$ the premise $\mathbb{E}[e_i\mid\mathcal{M}_j]\le1$
is exactly $\Pr(K_i\le k\mid\mathcal{M}_j)\le k/(\nCal+1)$, since $\ceil=(\nCal+1)/k$.
\end{proof}

\emph{Why the conditioning is per hypothesis.} A single global metadata $\sigma$-field $\mathcal{M}$
will not support the attack analysis of \cref{sec:attack}. $\mathbb{E}[e_i\mid\mathcal{M}]\le1$ does
\emph{not} imply $\mathbb{E}[e_i\mid\mathcal{M}']\le1$ for a richer $\mathcal{M}'$ --- take
$e=2\mathbf 1_A$ with $\Pr(A)=\tfrac12$: against the trivial $\sigma$-field the conditional
expectation is $1$, against $\sigma(A)$ it is $2$ on $A$ --- so an adversary who enriches the stream's
metadata anywhere would, on the global reading, void the premise even for groups it never touched.
With $\mathcal{M}_j$ the statement is local and survives: padding adds flows carrying the attacked
group's \emph{own} key, so it changes no other group's key, membership or arity, and adds no
hypothesis, leaving $T$ fixed. Provided the within-bucket order is a function of group metadata alone
(the deterministic hashed-key canonical order of \cref{sec:transfer}, with collisions broken on the
key itself), every true-null group's position and $\gamma_j$ are likewise fixed, so $\mathcal{M}_j$ is
literally the same $\sigma$-field before and after the attack. Under an arrival-ordered tie-break this
fails: a pad placed before the attacked group's first flow moves that group within its bucket and can
shift a true-null group's index, hence $\gamma_j$, hence $\mathcal{M}_j$. Creating a \emph{new} host
pair is likewise outside the padding claim --- it adds a hypothesis and changes $T$, which the
horizon-uniform $\gamma_j=1/T$ reads.

\emph{Scope.} The lemma certifies the \emph{e-value} controllers (e-LOND, online e-BH) only. The
p-value procedures carry their own conditions --- conditional super-uniformity for mFDR, and
independence with monotonicity for $\FDR$ --- which this argument does not supply, and which
\cref{sec:escapes} shows ADDIS violates on this stream. The premise
$\mathbb{E}[e_i\mid\mathcal{M}_j]\le1$ is marginal over the calibration draw and test score; it is
\emph{not} validity conditional on the realised calibration scores, which are deliberately excluded
from $\mathcal{M}_j$ (the Bates adjustment supplies that stronger, calibration-conditional guarantee at
a power cost and is not required here). $\mathcal{M}_j$ is also \emph{smaller} than a global
metadata $\sigma$-field, so \cref{assump:groupval} is the weaker of the two conditional premises,
while remaining stronger than unconditional $\mathbb{E}[e_i]\le1$.

\subsection{Feasibility (\cref{thm:family1,thm:family2}, \cref{cor:budget})}

\begin{proof}[\cref{thm:family1}]
With $\alphat = \alpha\gamma_t\,g(\mathcal{H}_{t-1})$, $g \le c_g(R_{t-1}+1)^d$, and
$R_{t-1}\le t-1$, we have $\alphat \le \alpha c_g\,\gamma_t t^d$. A rejection at $t$ needs
$\alphat \ge 1/\ceil$, hence $\gamma_t t^d \ge 1/(\alpha c_g\,\ceil)$. If $\gamma_t t^d \to 0$ then only
finitely many $t$ satisfy that inequality, since a sequence tending to zero exceeds a fixed positive
constant finitely often. A finite set of indices has a maximum, so there is a last feasible index
$t^\star$ (or none at all), and every $t>t^\star$ fails the necessary condition: the run is absorbed
from $t^\star+1$. No monotonicity of $\gamma$ is needed. Requiring permanence from the \emph{first}
infeasible step would be a stronger property --- false for non-monotone $\gamma$, since feasibility can
fail once and return --- and it is not what \cref{sec:feasibility} defines.
\end{proof}

\begin{proof}[\cref{thm:family2}]
Let $\alphat = w_0\gamma_t + \sum_j a_j\gamma_{t-\tau_j}$ with $w_0\ge0$ and $0\le a_j\le\alpha$ over
rejection times $\tau_j<t$. During a rejection-free run of length $\Delta$ the last rejection is at least
$\Delta$ before $t$, so every lag $t-\tau_j\ge\Delta$ and the lags are distinct. With
$\Gamma(\Delta)=\sum_{u\ge\Delta}\gamma_u$ we have $\gamma_{t-\tau_j}\le\Gamma(\Delta)$ termwise and
$\gamma_t\le\Gamma(t)\le\Gamma(\Delta)$, so, using $a_j\le\alpha$ and the distinctness of the lags,
$\alphat \le w_0\Gamma(\Delta) + \alpha\Gamma(\Delta) = (w_0+\alpha)\Gamma(\Delta)$; the first term is
where $w_0\ge0$ is used, since $\gamma_t\le\Gamma(\Delta)$ may only be multiplied through by a
non-negative $w_0$. Since $\Gamma$ is
non-increasing and $\sum\gamma<\infty$ forces $\Gamma(\Delta)\to0$, the quantity
$\Delta^\star=\min\{\Delta:(w_0+\alpha)\Gamma(\Delta)<1/\ceil\}$ is finite; $\Delta$ only grows during
a rejection-free run, so the state is absorbing. No monotonicity of $\gamma$ is used.
\end{proof}

\begin{proof}[\cref{cor:budget} and max-min optimality of $\gamma_t=1/T$]
Before any rejection, family~I has $\alphat=\alpha\gamma_t$ and family~II has $\alphat=w_0\gamma_t$;
write the cold-start coefficient $c_0\in\{w_0,\alpha\}$, so $\alpha_T=c_0\gamma_T$. Feasibility at $T$
requires $\alpha_T\ge 1/\ceil=k/(\nCal+1)$, i.e. $\nCal\ge k/\alpha_T-1=k/(c_0\gamma_T)-1$. Under
$\gamma_T=1/T$ this is $\nCal\ge kT/c_0-1$: $kT/w_0-1$ for a level-$w_0$ procedure such as LORD++,
and $kT/\alpha-1$ for LOND/e-LOND. For optimality, a cold-start rejection is possible at every
$t\le T$ iff $\min_{t\le T}\gamma_t \ge k/((\nCal+1)c_0)$; among sequences with
$\sum_{t\le T}\gamma_t\le 1$ the value $\min_{t\le T}\gamma_t$ is maximised, at $1/T$, by the
uniform allocation: $\min_{t\le T}\gamma_t\le T^{-1}\sum_{t\le T}\gamma_t\le1/T$, since a minimum is
at most an average, and $\gamma_t=1/T$ attains it. Constraining the mass over all $t$ rather than over
$t\le T$ changes nothing, since $\sum_t\gamma_t\le1$ implies $\sum_{t\le T}\gamma_t\le1$, and the
attaining sequence satisfies the global constraint too. Hence $\gamma_t=1/T$ maximises the feasible
horizon. That same criterion supplies the corollary's \emph{sufficiency}, which is otherwise easy to
miss: under $\gamma_t=1/T$ it reads $1/T\ge k/((\nCal+1)c_0)$, i.e.\ $\nCal\ge kT/c_0-1$ once more, so
this one inequality is not merely necessary at step $T$ but sufficient for feasibility at every
$t\le T$. The two directions coincide here only because a cold-start uniform level is constant over
the horizon; for a general $\gamma$ the necessary condition at $T$ is strictly weaker than
feasibility throughout, which is why the corollary states the equivalence for $\gamma_t=1/T$ alone.
\end{proof}

\subsection{The UAI 2026 procedures (\cref{prop:donation,prop:closure})}

\begin{proof}[\Cref{prop:donation}]
Donation e-LOND sets
$\alphat = \delta\gamma_t(\lvert R_{t-1}\rvert+1) / \bigl(1 - (\delta(\lvert R_{t-1}\rvert+1)\bar W_t \wedge 1)\bigr)$
with
$\bar W_t = \sum_{i\in R_{t-1}} \gamma_i\bigl((E_i - \tfrac{1}{\delta\gamma_i(\lvert R_{t-1}\rvert+1)}) \wedge 1\bigr)
+ \sum_{i\notin R_{t-1}} \gamma_i (E_i \wedge 1)$.
with both sums running over $i<t$. Every summand is $\gamma_i$ times a quantity capped at $1$, so
$\bar W_t \le \sum_{i<t} \gamma_i \le 1$ \emph{always}. On a rejection-free prefix
$R_{t-1}=\emptyset$, so for $\delta<1$ the denominator is at least $1-\delta>0$ and
$\alphat \le \delta\gamma_t/(1-\delta)$, which is \cref{thm:family1} with
$c_g = \delta/[\alpha(1-\delta)]$ and $d = 0$: \cref{thm:family1} writes the level as
$\alpha\gamma_t g(\mathcal{H}_{t-1})$, so the constant must absorb the ratio $\delta/\alpha$, and it
reduces to $1/(1-\delta)$ exactly in the $\delta=\alpha$ setting we run. $d=0$ is the right
instantiation, not merely a valid one: on a rejection-free prefix $R_{t-1}=0$, so the displayed bound
$g\le c_g$ holds uniformly in $t$ --- $g$ itself rises with $\bar W_t$ and is not constant --- and
\cref{thm:family1}'s conclusion then needs only $\gamma_t\to0$, which $\sum_t\gamma_t\le1$ gives.
Taking $d=1$ instead would demand $t\gamma_t\to0$, which summability does \emph{not} imply
($\gamma_{2^j}=2^{-j}$, zero elsewhere, sums to one with $t\gamma_t=1$ infinitely often).

\emph{Why $\delta<1$ is required.} At $\delta=1$ the denominator is $1-\bar W_t$, which the donated
wealth drives to zero, and the conclusion fails outright. Take $\delta=1$, $\gamma_t=2^{-t}$ and
$E_t=1$ for every $t$. Then $\bar W_t=\sum_{i<t}\gamma_i(E_i\wedge1)=\sum_{i<t}\gamma_i=1-2^{-(t-1)}$
and $\alphat=\gamma_t/(1-\bar W_t)=2^{-t}/2^{-(t-1)}=\tfrac12$ for \emph{every} $t$. The run is
rejection-free because rejection needs $E_t\ge1/\alphat=2$ and $E_t=1$ --- this is why the evidence is
pinned at exactly $1$ rather than merely bounded below by it --- yet the offered level never decays, so
with any $\ceil\ge2$ every step stays evidence-feasible and the run is never absorbed. The proposition
therefore quantifies over $\delta\in(0,1)$; our runs use $\delta=\alpha=0.05$, where
$1/(1-\delta)=1.05$.

The restriction to the rejection-free prefix is not cosmetic. With $R = \lvert R_{t-1}\rvert$ the
denominator is $1 - (\delta(R+1)\bar W_t \wedge 1)$, which vanishes once $\delta(R+1)\bar W_t \ge 1$;
since $\bar W_t$ can approach $1$, that is reachable at $R+1 \ge 1/\delta$. The bound is therefore a
cold-start statement, which is exactly the state \cref{cor:budget} concerns. It is also sharp: taking
$E_i \ge 1$ for every $i$ makes $\bar W_t = \sum\gamma_i \to 1$ and the bound tight.
\end{proof}

\begin{proof}[\Cref{prop:closure}]
Closed e-LOND sets
$\alphat = \min_{S\subseteq[t-1]:\,D_t(S)>0} \delta\gamma_{\lvert S\rvert+1}(\lvert R_{t-1}\rvert+1)/D_t(S)$
with $D_t(S) = 1 + \lvert S\cap R_{t-1}\rvert - \delta E_S(\lvert R_{t-1}\rvert+1)$, where $E_S$
aggregates the evidence on $S$. Because $\alphat$ is a \emph{minimum} over admissible subsets, any
single admissible $S$ furnishes an upper bound.

Take $S = \{i<t : E_i = 0\}$, of size $Z_t$. A hypothesis carrying zero evidence is never rejected, so
$S\cap R_{t-1} = \emptyset$; and $E_S = 0$ because the aggregator is zero on all-zero input,
including the empty subset that arises when $Z_t=0$. Hence
$D_t(S) = 1 > 0$, $S$ is admissible, and
$\alphat \le \delta\gamma_{Z_t+1}(\lvert R_{t-1}\rvert+1)$.

This is the decay of \cref{thm:family1} with the time index rescaled from $t$ to $Z_t+1$. The bound
itself is unconditional and pathwise; the \emph{conclusion} that the state is absorbing is not, and
\cref{cor:closureabsorb} supplies its hypothesis.
\end{proof}

\begin{proof}[\Cref{cor:closureabsorb}]
By \cref{prop:closure}, $\alphat \le \delta\gamma_{Z_t+1}(\lvert R_{t-1}\rvert+1)$ along any run. On a
rejection-free run $\lvert R_{t-1}\rvert = 0$, so $\alphat \le \delta\gamma_{Z_t+1}$. If
$\gamma_{Z_t+1}\to0$ then $\alphat\to0$, and since the evidence is bounded above by $\ceil$
(\cref{sec:feasibility}) the offered level falls below $1/\ceil$ after finitely many steps and no
rejection is possible thereafter: the state is absorbing, which is \cref{thm:family1}'s conclusion.

\emph{Form (i), exact and finite-horizon.} Along the rejection-free run $\lvert R_{t-1}\rvert=0$, so
$\alphat\le\delta\gamma_{Z_t+1}$. Our $\gamma_j\propto j^{-1.6}$ is decreasing, so no rejection is
possible once $\gamma_{Z_t+1} < 1/(\delta\ceil)$, i.e.\ once $Z_t \ge H-1$ for
$H = \min\{j : \gamma_j < 1/(\delta\ceil)\}$, which is exactly \cref{thm:family1}'s index horizon.
Write $P_T = \#\{i\le T : E_i>0\}$. Of the $t-1$ hypotheses before step $t$, at most $P_T$ carry
positive evidence, so $Z_t \ge t-1-P_T$ for every $t\le T$ --- a counting statement, with no
assumption on the stream. Hence $Z_t \ge H-1$ as soon as $t \ge H + P_T$, and
\[
  t_{\text{closed}} \ :=\ \inf\{t : Z_t \ge H-1\} \ \le\ H + P_T .
\]
The closure therefore reaches the absorbing state within an \emph{additive} $P_T$ steps of the horizon
\cref{thm:family1} gives, whenever $H+P_T\le T$. Nothing is asymptotic and no density is assumed. For a
non-monotone $\gamma$ the same conclusion follows without the decreasing step, because $\alphat$ is a
\emph{minimum} over admissible subsets: any $H-1$ zero-evidence indices furnish an admissible $S$ and
hence the same bound.

\emph{Form (ii), asymptotic.} Suppose $\liminf_t Z_t/t = \rho > 0$ and fix any $\rho' < \rho$. By the
definition of the $\liminf$ there is a $t_1$ with $Z_t \ge \rho' t$ for all $t \ge t_1$; so
$Z_t\to\infty$ and, $\gamma_j\propto j^{-1.6}$ being decreasing to $0$,
$\gamma_{Z_t+1} \le \gamma_{\lceil\rho' t\rceil}\to0$. Substituting into \cref{thm:family1}'s
discovery-horizon computation replaces the index $t$ by $\rho' t$: the effective spending index is
rescaled by $\rho'$, so absorption is reached once both $t\ge t_1$ and $\rho' t\ge H-1$, i.e.\ by
$\max\{t_1,\lceil(H-1)/\rho'\rceil\}$. The onset $t_1$ cannot be dropped: a positive $\liminf$
constrains only the eventual density and permits an arbitrarily long initial transient. The
\emph{time} to absorption is therefore stretched by $1/\rho'$, not shortened --- a sparser
zero-evidence stream buys the closure more steps, not fewer. This holds for every $\rho'<\rho$ and \emph{not}, from this argument, at $\rho$ itself:
the liminf gives eventual domination below its value, not at it.

On our streams form \emph{(i)} is the operative one and needs no density at all. The number of
hypotheses carrying \emph{any} evidence is $P_T \in \{0,\dots,152\}$ over the ten window--seed cells,
at most $\hat p = 4.82\times10^{-3}$ of $T$, attained at position $0.85$, and $P_T=0$ at $0.70$ seed
$1$; so the horizon offset is at most $152$ hypotheses and the cap is attained (ratio $1.000000$ over
every configuration measured). The argument is specific to \emph{sparse} evidence: a stream in which
most hypotheses carry positive evidence has $P_T$ close to $T$, $Z_t$ small, and this bound says
little --- which is the regime the closure was designed for, and not a security stream.
\end{proof}

\subsection{Symmetric padding impossibility (\cref{thm:padding})}

We give a full proof (Route A) and an independent lemma (Route B) that needs no domination result.
Route B's scope is narrow and worth stating exactly: it shows that a symmetric family maps any
permutation of a mean-one two-point vector to at most $1$. That yields the impossibility only for
families whose \emph{unpadded} value on that vector already reaches $\theta$ --- the arithmetic mean
among them --- and not for the class. A symmetric valid family such as
$F_N(e)=\mathrm{mean}(e)\,\mathbf 1\{\text{not all }e_i\text{ equal}\}$ attains $\theta$ elsewhere yet
returns $0$ on Route B's unpadded vector, so Route B produces no violation for it and Route A must.
The class-wide statement and the quantitative bound are Route A's alone.

\begin{proof}[Route A (essential domination)]
Vovk and Wang's Proposition~3.1~\cite{vovk2021evalues} states that the arithmetic mean $M_K$
\emph{essentially} dominates every symmetric e-merging function: $F(e)>1\Rightarrow M_K(e)\ge F(e)$.
Ordinary domination is false ($\kappa+(1-\kappa)M_K$, for any $\kappa\in(0,1)$, is a symmetric valid counterexample), so we
use it through a case split: for finite $x$, either $F_N(x)\le1$, or $F_N(x)>1$ and then
$M_K(x)\ge F_N(x)$; in both cases $F_N(x)\le\max\!\big(1,\mathrm{mean}(x)\big)$.

We need that inequality at $(x,0^r)$, a boundary point of $[0,\infty)^N$, so we prove it directly
rather than rely on the domain of the cited proposition. Let $y\in[0,\infty)^N$, $\bar y=\mathrm{mean}(y)$,
and let $\Pi$ be a uniformly random permutation of $y$; each coordinate of $\Pi$ has mean $\bar y$, and
$F_N(\Pi)\equiv F_N(y)$ by symmetry. If $\bar y\le1$ then $\Pi$ is an admissible input, $F_N(\Pi)$ is a
constant e-variable, so $F_N(y)\le1$. If $\bar y>1$, take $E=\Pi$ with probability $1/\bar y$ and
$E=0$ otherwise; every coordinate of $E$ then has mean exactly $1$, so validity gives
$F_N(y)/\bar y+(1-1/\bar y)F_N(0)\le1$, and $F_N(0)\ge0$ yields $F_N(y)\le\bar y$. Either way
$F_N(y)\le\max(1,\bar y)$ on all of $[0,\infty)^N$, zeros included, which is what the padding step uses.
This also removes any need for monotonicity or measurability assumptions on $F_N$. Padding
$x\in[0,\infty)^m$ with $r$ zeros gives $\mathrm{mean}(x,0^r)=\big(\sum_i x_i\big)/(m+r)$, so
\[
F_{m+r}(x,0^r)\ \le\ \max\!\Big(1,\ \tfrac{\sum_i x_i}{m+r}\Big)\ <\ \theta
\ \text{ once }\ r>\tfrac{\sum_i x_i}{\theta}-m .
\]
Attainment supplies an $x$ with $F_m(x)\ge\theta$ to which this applies, so no attaining symmetric
family is $\theta$-padding-robust for any $\theta>1$.
\end{proof}

\begin{proof}[Route B (a negatively-dependent special case)]
Fix $\theta>1$ and $m$. Draw a uniformly random size-$m$ subset $S\subseteq[N]$ and set
$e_i=(N/m)\mathbf{1}\{i\in S\}$. Then $\mathbb{E}[e_i]=(N/m)\Pr(i\in S)=1$, so each $e_i$ is an
e-variable, and the vector is a permutation of $\big((N/m)\mathbf{1}_m,0^{N-m}\big)$. A symmetric
$F_N$ maps every realisation to the same constant, and a constant e-variable is $\le1$. Hence
$F_N\big((N/m)\mathbf{1}_m,0^{N-m}\big)\le1<\theta$ for all $N\ge m$: appending $N-m$ zeros to the
$m$ firing positions drives any symmetric rule below $\theta$.
\end{proof}

\emph{Tightness and attainment.} The bound is tight: the mean itself has $F_{m+r}(x,0^r)=(\sum
x)/(m+r)\ge\theta$ exactly while $r\le(\sum x)/\theta-m$. Attainment (\cref{sec:attack}) is necessary,
not cosmetic: the constant rule $F\equiv1$ is symmetric, valid and padding-invariant but never
reaches any $\theta>1$, satisfying $\theta$-padding-robustness vacuously; since rejection needs
evidence $\ge1/\alphat>1$, restricting to attaining families excludes exactly the rules that
cannot fire.

\subsection{Feasibility ratio as a dilution threshold (\cref{cor:dilution})}

\begin{proof}[\cref{cor:dilution}]
Under $\gamma_t=1/T$ the level is $\alphat=(\alpha/T)(R_{t-1}+1)$, so on a rejection-free prefix
$\alphat=\alpha/T$ and \eqref{eq:reject} rejects $G_t$ exactly when $\Ev(G_t)\ge T/\alpha$; the rule
is inclusive. \emph{Sufficiency.} Let $c>\rho$ be an integer; we show by induction on $t$ that no
rejection occurs. Suppose none has occurred before $t$, so $R_{t-1}=0$ and the level is $\alpha/T$.
If $G_t$ is not an attacker episode, $\Ev(G_t)<T/\alpha$ by hypothesis. If it is, with unpadded sum
$S_t$ and arity $m_t$, each of its flow e-values is at most $\ceil$ by \eqref{eq:ceiling}, so
$S_t\le\ceil m_t$; its padded arity is the integer $cm_t$ and the padded mean is
$S_t/(cm_t)\le\ceil/c<\ceil/\rho=T/\alpha$, using $c>\rho$ and $\rho=\ceil\alpha/T$. In either case
\eqref{eq:reject} fails at $t$ and the induction closes. \emph{Necessity.} Let $c\le\rho$ be an
integer and let some attacker episode $G_u$ attain the ceiling, $S_u=\ceil m_u$. If the attacked
trajectory were rejection-free, $G_u$ would be tested at $\alpha/T$ with padded mean
$\ceil m_u/(cm_u)=\ceil/c\ge\ceil/\rho=T/\alpha$ and rejected by the inclusive rule, a
contradiction. Hence no integer $c\le\rho$ suffices for every stream satisfying the hypotheses,
whereas the integer $\lfloor\rho\rfloor+1>\rho$ does. Two remarks. The hypothesis on the other
hypotheses cannot be dropped: a benign group rejected at the cold-start level raises $R$ and with it
every later level, and the attacker's pads are then priced above $\alpha/T$. And the threshold is
uniform over streams: for one given stream with at least one attacker episode, the same two
arguments, with $\ceil$ replaced by $\max_u\Ev(G_u)$ over the attacker's episodes, show that the
smallest sufficient integer is
$\lfloor c_{\mathrm{crit}}\rfloor+1$ with $c_{\mathrm{crit}}=\max_u\Ev(G_u)\,\alpha/T\le\rho$, and
$c_{\mathrm{crit}}=\rho$ exactly when some attacker episode attains $\ceil$. A real-valued $c$ would
make $(c-1)m_t$ non-integral; the statement is therefore made for integers, which is what an attacker
can append.
\end{proof}

\subsection{Reach and front-load (\cref{thm:reach,thm:frontload})}

Let $w=(w_1,w_2,\dots)$ be precommitted with $w_i\ge0$ and $\sum_i w_i\le1$, let
$F(e)=\sum_i w_i e_i$ with each $e_i\in\{0,\ceil\}$, and write $\beta=1/(\alphat\ceil)$ for the
weight mass a detection needs: a lone flow firing at position $i$ gives $F=w_i\ceil$, which clears
$1/\alphat$ iff $w_i\ge\beta$.

\begin{proof}[\cref{thm:reach}]
The positions at which a lone attack flow fires are $\{i:w_i\ge\beta\}$. Since $\sum_i w_i\le1$,
at most $1/\beta$ indices can carry mass $\ge\beta$, so
$P=|\{i:w_i\ge\beta\}|\le 1/\beta=\alphat\ceil$, the quantity that governs \cref{cor:budget}.
\end{proof}

\begin{proof}[\cref{thm:frontload}]
Prepending $L$ zero-evidence flows moves every attack flow to a position $>L$. Every $e_i\le\ceil$, so
$F(e)\le\ceil\sum_{i>L}w_i$, and detection requires $\sum_{i>L}w_i\ge\beta$; it is impossible once
$L\ge L^\star(w)=\min\{L:\sum_{i>L}w_i<\beta\}$. Summability gives $\sum_{i>L}w_i\to0$, so the set is
non-empty and $L^\star$ is finite for every summable $w$.

\emph{Sufficient, and exact only under a full-firing episode.} The displayed bound is attained only when
every position $>L$ fires, so $L^\star$ is the exact minimum in that case and an \emph{upper} bound
otherwise. The gap is real: for $w=(0.4,0.3,0.3,0,\dots)$ and $\beta=0.5$ the tail first falls below
$\beta$ at $L=2$, yet an episode carrying a \emph{lone} firing flow is already undetectable at $L=0$,
since $\max_i w_i=0.4<\beta$ and \cref{thm:reach} gives $P=0$. We therefore quote $L^\star$ as an
attacker's upper bound throughout.
\end{proof}

Consequently a precommitted scheme cannot simultaneously be padding-invariant (weights fixed by
absolute position), have usable power (some $w_i\ge\beta$), and make ordering robustness unbounded
($L^\star=\infty$). Concentrating weight early gives $L^\star=1$ but $P=1$; spreading it over more
than $1/\beta$ positions drives every $w_i<\beta$, so $P=0$.

\subsection{ADDIS state budget (\cref{thm:addis}) and the coincidence}

\begin{proof}[\cref{thm:addis}]
ADDIS offers $\alphat=\min\{\lambda,(\tau-\lambda)[W\gamma_{S^t-C_{0+}}+\cdots]\}$. The theorem
assumes $\lambda\ge k/(\nCal+1)$, i.e.\ the cap sits at or above the conformal floor, so that the
driving term and not the cap is what the attack has to push below $1/\ceil$; if the cap were already
below the floor no target could be rejected at $D=0$ and the budget would be zero rather than the
displayed value. With $W=w_0$
before the first rejection and $W=\alpha R$ after $R$ of them, and $0$-indexed
$\gamma_j=(j+1)^{-1.6}/\zeta(1.6)$. Precursors injected at index $0$ hold the spending index at
$S^t-C_{0+}=D$, so the driving level is $(\tau-\lambda)W\gamma_D$. The collapse of ADDIS's reward sum to
the single term $W\gamma_D$ is exact in the pre-rejection state, where $W=w_0$ and there is one active
term. After $R\ge1$ rejections each reward term carries its own selected-non-candidate lag
$\ell_j\ge0$, sitting at index $D-\ell_j$; the coefficients sum to exactly
$w_0+(\alpha-w_0)+\alpha(R-1)=\alpha R$, but $\gamma$ is decreasing, so
$\sum_j a_j\gamma_{D-\ell_j}\ge\alpha R\,\gamma_D$. The drive is therefore \emph{larger} than the
collapsed expression, feasibility survives further, and the displayed value at $W=\alpha R$ is a
\emph{lower} bound on $B^\star$ --- equality only when every $\ell_j=0$. Numerically, at
$\lambda\ceil=(\tau-\lambda)\ceil=4.5\times10^{5}$, $R=1$ and a single lag of $100$, the expression
returns $313$ while the true budget is $373$. The attack is mounted in the pre-rejection state, since the precursors precede every target. A rejection is possible only
while $(\tau-\lambda)W\gamma_D\ge1/\ceil$, i.e.
$(D+1)^{-1.6}/\zeta(1.6)\ge k/\big((\tau-\lambda)W(\nCal+1)\big)$, i.e.
$D+1\le[(\tau-\lambda)W(\nCal+1)/(k\zeta(1.6))]^{1/1.6}$. The smallest $D$ violating this is
$B^\star$. Each precursor is a group carrying exactly one firing flow, so its evidence is $\ceil/m$
and its p-value $m/\ceil$; landing in $(\lambda,\tau]$ needs $m\in(\lambda\ceil,\tau\ceil]$, whose
smallest integer is $m^\star=\lfloor\lambda(\nCal+1)/k\rfloor+1$ (the interval is non-empty because
$(\tau-\lambda)\ceil\gg1$). The total is $B^\star m^\star$, scaling as
$\lambda\nCal\,[(\tau-\lambda)w_0\nCal]^{0.625}/k^{1.625}$.
\end{proof}

\emph{Coincidence.} When ADDIS's cap binds, $\alphat=\lambda$, and the minimal suppressing pad
raises the padded p-value just past $\lambda$. Because the p-value moves in steps of one conformal
floor $1/(\nCal+1)\ll\tau-\lambda$, it lands in $(\lambda,\tau]$ and advances the index:
cap-saturation \emph{implies} landing. The converse is not forced, but on this stream every
unsaturated detection has a level four orders of magnitude below $\lambda$, so none can land, and
the observed equality (all $101$ saturated detections land) is the generic case rather than a
coincidence of this window.

\subsection{Restart FDR composition}

\begin{proof}
Restart runs $n$ controllers on disjoint epochs, each controlling its own $\FDR$ at $q$. The
deployment-wide proportion $\FDP=\sum_i V_i/(\sum_i R_i\vee1)$ is a weighted average of the
per-epoch $\FDP_i$ with data-dependent weights $R_i/\sum_j R_j$, so per-epoch control does not
bound $\mathbb{E}[\FDP]$. Indeed, suppose each epoch makes \emph{one false rejection with
probability $q$ and no rejection otherwise}. Each epoch then has
$\FDR_i=\mathbb{E}[V_i/(R_i\vee1)]=q$, yet whenever any epoch rejects every pooled rejection is
false, so the deployment-wide $\FDP$ is $1$ on that event and $0$ otherwise, and with independent
epochs $\FDR=\Pr(\text{some epoch rejects})=1-(1-q)^n$. (The neighbouring construction --- each epoch
\emph{always} rejecting once, false with probability $q$ --- instead gives pooled $\FDR=q$, so
non-composition requires the all-or-nothing epoch above and does not follow from per-epoch control
alone.) In contrast, mFDR pools: if
$\mathbb{E}[V_i]\le q\,\mathbb{E}[R_i]$ for every $i$, the mediant inequality gives
$\sum_i\mathbb{E}[V_i]/\sum_i\mathbb{E}[R_i]\le q$.

This non-composition is specific to \emph{budget reset} --- every epoch spends its own $q$, total
$nq$. A \emph{preallocated} restart --- pre-committed epoch budgets $q_1,\dots,q_n$ with
$\sum_i q_i\le q$, fixed before deployment --- instead \emph{preserves} the deployment-wide guarantee.
Since $\sum_j R_j\ge R_i$, the pooled proportion is dominated termwise,
\[
\FDP=\frac{\sum_i V_i}{\sum_j R_j\vee1}\le\sum_i\frac{V_i}{R_i\vee1}=\sum_i\FDP_i
\]
(both sides $0$ when no epoch rejects), so by linearity of expectation
$\FDR\le\sum_i\FDR_i\le\sum_i q_i\le q$ --- no independence across epochs is required (independence
enters only the $1-(1-q)^n$ counterexample above). Each e-LOND epoch delivers $\FDR_i\le q_i$ under the
same group e-validity premise (\cref{assump:groupval}); the shared calibration set only induces
between-epoch dependence, which this bound tolerates, provided the epoch boundaries and the $q_i$ are
not data-dependent. Restart therefore does not simply void the deployment-wide guarantee: it forces a
\emph{choice} --- budget reset recovers more power but controls only a per-epoch $\FDR$, whereas
preallocation keeps deployment-wide $\FDR\le q$ at a power cost (\cref{apptab:restart}).
\end{proof}
